\section{El tercer paradigma, Online Learning: por qué podría ser más fundamental que los robots, los modelos del mundo y la multimodalidad}

Lo anterior trató el panorama competitivo actual; esta sección entra en los paradigmas técnicos. Guangmi (广密) considera que el scaling del preentrenamiento está por terminar y que RL/Post-training todavía tiene margen, pero que lo que realmente podría cambiar el panorama en la siguiente etapa es el Online Learning o Continuous Learning. Lo central aquí no se reduce a que «el modelo tenga un contexto más largo», sino a si el modelo puede aprender de forma continua a partir de la interacción, la retroalimentación y las tareas reales.

\begin{figure}[H]
\centering
\includegraphics[width=0.96\textwidth]{scaling-paradigm-shift.png}
\caption{Tres paradigmas de aprendizaje: preentrenamiento, RL/Post-training, Memory y Online Learning. Diagrama conceptual propio, elaborado a partir de la conversación de 00:36:01--00:41:01.}
\end{figure}

\begin{knowledgebox}{Lectura de la figura: de la compresión estática al aprendizaje por interacción}
El preentrenamiento equivale a comprimir de una sola vez datos de internet y datos especializados a gran escala; RL/Post-training moldea el modelo con retroalimentación humana, procesos de expertos y resultados de tareas; Memory y el contexto largo permiten que el modelo recuerde más historial del usuario y de las tareas; Online Learning va más allá y exige que el modelo se actualice realmente durante la inferencia y la interacción, de modo que entienda mejor al usuario cuanto más se usa y haga mejor las tareas cuanto más las realiza.
\end{knowledgebox}

\subsection{Preentrenamiento: un compresor gigantesco}

El preentrenamiento (Pre-training) consiste en entrenar un modelo con cantidades masivas de texto, código y datos multimodales para que aprenda a predecir el siguiente token o la siguiente estructura. Guangmi enfatiza que los modelos actuales siguen siendo compresores gigantescos: las tareas que no están en la distribución de los datos suelen no funcionar, y el modelo solo es más confiable cuando ha comprimido datos similares.

Esta afirmación es fundamental. Explica por qué los modelos ya superan a la mayoría de las personas en conocimiento y, aun así, siguen fallando en flujos de trabajo reales. El trabajo real no consiste solo en responder preguntas, sino en operar dentro de un CRM, sistemas financieros, software de diseño, procesos hospitalarios, ventas en campo, permisos empresariales y cadenas de herramientas concretas. Si el modelo no ha visto suficientes procesos reales, solo puede recurrir a conjeturas por generalización.

\begin{importantbox}{La perspectiva del compresor}
El modelo no «entiende» una tarea de la nada, sino que comprime los patrones de los datos de entrenamiento en parámetros y representaciones. Si una tarea nunca ha entrado en la distribución de los datos, al modelo le resulta difícil completarla de forma estable; para que un Agent sea confiable, hay que convertir los procesos de trabajo reales en datos de los que se pueda aprender.
\end{importantbox}

\subsection{RL/Post-training: datos de expertos y retroalimentación de tareas}

Antes se explicó el preentrenamiento como una compresión a gran escala; esta subsección examina «cómo se moldea el modelo después de la compresión». Si el preentrenamiento permite que el modelo aprenda la gran cantidad de patrones que ya existen en el mundo, RL/Post-training se ocupa más de cómo el modelo se vuelve útil en tareas concretas, ante las preferencias del usuario y en los procesos de expertos.

RL es Reinforcement Learning, aprendizaje por refuerzo. No se reduce a las palabras «modelo de recompensa», sino que es una forma de entrenamiento en la que el sistema ajusta su comportamiento según la retroalimentación de los resultados. En el contexto de los modelos grandes, Post-training suele incluir ajuste por instrucciones, retroalimentación de preferencias humanas, trayectorias de expertos, evaluación automática, retroalimentación sobre el uso de herramientas y calificación de resultados de tareas. En el programa se compara este tipo de datos con las «nuevas energías»: son útiles, pero su volumen total es reducido y su recolección es costosa.

Para que un Coding Agent mejore, se necesitan código de alta calidad, issues reales, pruebas, retroalimentación de PR y procesos de ingeniería de software de largo plazo; para que un Agent médico mejore, se necesitan los procesos de diagnóstico de los médicos, expedientes clínicos, estudios, medicación, seguimiento y límites de responsabilidad; para que un Agent de análisis de inversiones mejore, se necesitan los flujos de trabajo de los analistas, las fuentes de datos, los supuestos de los modelos y la revisión de errores. Ninguna de estas categorías es texto web común: todas son datos de procesos de expertos.

\begin{knowledgebox}{Términos clave: vocabulario de RL/Post-training}
\begin{tabularx}{\textwidth}{p{0.22\textwidth}p{0.30\textwidth}X}
\toprule
Término & Mecanismo central & Papel en este episodio \\
\midrule
RL & Optimiza la política con recompensas o retroalimentación de resultados & Hace que el modelo no solo imite texto, sino que se optimice hacia el éxito en la tarea. \\
Post-training & Alineación, instrucciones, preferencias y entrenamiento por tareas posteriores al preentrenamiento & Determina si el modelo se ajusta a los objetivos del producto. \\
Trayectorias de expertos & Registro del proceso con el que expertos de primer nivel completan una tarea & Permite que el modelo aprenda flujos de trabajo reales, no solo respuestas. \\
Credit assignment & Atribuir el éxito o el fracaso a los pasos intermedios & En la sección de Robotics se usa para explicar si cada paso favorece el éxito de la tarea. \\
Value function & Función que estima el valor futuro del estado o la acción actuales & Ayuda al sistema a saber qué pasos están más cerca del éxito. \\
\bottomrule
\end{tabularx}
\end{knowledgebox}

\subsection{Online Learning: aprender mientras se interactúa}

En este episodio, Online Learning se refiere a que el modelo aprenda de manera autónoma y continua, y que absorba la experiencia del usuario y de las tareas a partir de la interacción. A diferencia del preentrenamiento estático, exige que el modelo, a medida que se usa, entienda mejor al usuario y la tarea y actúe con más iniciativa. El escenario de producto que plantea Guangmi es este: hoy ChatGPT tiene que esperar el prompt del usuario; en el futuro, la IA podría reflexionar, prepararse y establecer asociaciones en segundo plano, de modo que, cuando el usuario vuelva a preguntar, ya tenga una comprensión más profunda.

Si esta capacidad se materializa, cambiará la forma de los productos. La IA dejará de ser una herramienta que razona desde cero cada vez y pasará a ser un compañero de largo plazo que crece junto con el usuario; dejará de limitarse a responder la pregunta actual y recordará los objetivos, preferencias, proyectos, contexto y errores pasados del usuario; dejará de esperar órdenes de forma pasiva y se preparará y emitirá recordatorios por iniciativa propia.

\begin{importantbox}{Definición de Online Learning}
En estos apuntes, Online Learning designa la capacidad de un modelo o sistema de IA de aprender de la interacción continua, de los resultados de las tareas y de la retroalimentación del usuario, y de aplicar lo aprendido de forma estable a su comportamiento futuro. No se trata de guardar el historial de conversaciones, sino de convertir la experiencia en actualizaciones de capacidades, preferencias o estrategias.
\end{importantbox}

\begin{warningbox}{Memory no equivale a Online Learning}
Memory puede limitarse a almacenar información del usuario; el contexto largo puede limitarse a meter más historial en el prompt. Online Learning exige que el sistema aprenda de ese historial mejoras de comportamiento generalizables. Saber «recordar hechos» no equivale a «volverse más inteligente cuanto más se usa».
\end{warningbox}

\subsection{La metáfora energética: petróleo, nuevas energías y fusión nuclear}

Esta subsección usa una metáfora energética para condensar las tres etapas de aprendizaje anteriores en un marco comparativo fácil de recordar. El valor de la metáfora no está en predecir con precisión la ruta tecnológica, sino en ayudar al lector a ver a la vez las diferencias de restricciones entre tres tipos de recursos: los datos existentes, los datos de expertos y el aprendizaje continuo.

Guangmi usa tres metáforas energéticas para explicar las diferencias entre paradigmas: los datos de preentrenamiento son como el petróleo, abundantes pero finitos; los datos de expertos para RL son como las nuevas energías, útiles pero escasos en total; Online Learning es como la fusión nuclear, que todavía no logra un avance real, pero que, si lo logra, tendrá una densidad energética y un techo completamente distintos.

\begin{figure}[H]
\centering
\includegraphics[width=0.96\textwidth]{energy-metaphor.png}
\caption{Metáfora energética: el petróleo, las nuevas energías y la fusión nuclear corresponden a tres tipos de datos y de formas de aprendizaje. Diagrama conceptual propio, elaborado a partir de la conversación de 00:41:01--00:43:05.}
\end{figure}

\begin{knowledgebox}{Lectura de la figura: los límites de la metáfora}
El petróleo subraya que los datos existentes son finitos; las nuevas energías, que los datos de expertos son útiles pero difíciles de recolectar; la fusión nuclear, el techo potencial de Online Learning. La metáfora no implica que la tecnología vaya a repetir necesariamente la historia de la energía; solo nos ayuda a distinguir las restricciones de tres tipos de datos y mecanismos de aprendizaje.
\end{knowledgebox}

\subsection{Por qué los robots, los modelos del mundo y la multimodalidad podrían ser «falsos problemas»}

La afirmación más polémica del programa es esta: muchos de los temas de robots, modelos del mundo y multimodalidad podrían ser falsos problemas, y Online Learning sería el verdadero problema. Aquí «falso problema» no significa que estas líneas no sean importantes, sino que podrían depender de una generalización y un aprendizaje continuo más fundamentales. Sin generalización, los robots acabarían como la conducción autónoma: cada escenario de cola larga requiere recolectar datos, etiquetarlos y ajustar reglas, y se cae por mucho tiempo en la situación de «tanta inteligencia como trabajo humano haya».

Con la multimodalidad y los modelos del mundo ocurre algo parecido. Que un modelo pueda ver video, generar imágenes y predecir procesos físicos no significa que pueda aprender de forma continua en entornos reales, transferir experiencia y resolver tareas nuevas. Lo verdaderamente difícil es convertir la experiencia en capacidades reutilizables. Si Online Learning logra un avance, los robots, los modelos del mundo y la multimodalidad podrían obtener una generalización mucho mayor; si no, estas líneas seguirán frenadas por el volumen de datos, el costo del etiquetado y la complejidad de los entornos reales.

\subsection{Resumen de la sección}

Lo central del tercer paradigma no es un término nuevo, sino el paso de la compresión estática al aprendizaje continuo por interacción. El preentrenamiento resuelve la compresión de conocimiento a gran escala; RL/Post-training resuelve la retroalimentación de tareas y los procesos de expertos; Online Learning intenta resolver el problema de «volverse más inteligente cuanto más se usa» y de «generalizar a partir de entornos reales». Su importancia radica en que podría alterar el equilibrio actual, en el que los tres grandes actores se alternan el liderazgo.
